\documentclass{nfr}
\geometry{margin=0.8in}

\title{Supporting Statement: Postdoctoral Research Associate in the Economics of Technological Change}
\author{Zhongli Wang}

\begin{document}

\maketitle

I am applying for the Postdoctoral Research Associate position in the Economics of Technological Change (vacancy 189392). I completed my PhD in Economics at Bielefeld University and Universit\`a Cattolica del Sacro Cuore in 2026, successfully defending my dissertation on 21 September. My research combines complexity economics, industrial organisation, quantitative modelling, and artificial intelligence. The Oxford Martin Programme on Forecasting Technological Change is therefore a close match for both my substantive interest in how technologies develop and diffuse and my methodological experience with computational models, empirical calibration, machine learning, and large language models (LLMs). I am particularly motivated by the programme's plan to connect a comprehensive evidence base on technological progress with testable questions about prediction, diffusion, and the feedback between falling costs and adoption.

\section*{Essential selection criteria}

\subsection*{PhD in a relevant discipline}

I hold a PhD in Economics completed through the Economic Policy in Complex Environments (EPOC) Marie Sk\l odowska-Curie doctoral programme. My dissertation, \textit{Strategic Behavior and Economic Dynamics in ICT Markets: An Agent-Based Framework}, consists of three single-authored studies of data governance, capacity investment, and parameter estimation in agent-based economic models. This work required me to define research questions independently, develop theory and computational methods, build and analyse models, and turn results into publication-ready papers. It provides a direct foundation for independent research on the economics of technological change.

\subsection*{Complex systems science and forecasting methods}

My doctoral training is centred on complex adaptive economic systems. I study markets in which heterogeneous consumers, competing firms, network effects, capacity constraints, regulation, and feedback mechanisms jointly produce aggregate outcomes. In my job-market paper, firms adjust investment in service capacity while consumer adoption responds to service quality, congestion, prices, and peer influence. The resulting dynamics are nonlinear and path-dependent, so they cannot be understood from a representative-agent or static perspective alone. This experience is directly relevant to analysing feedback between technology costs, performance, adoption, and diffusion.

I also work on the empirical estimation of agent-based models. My current research compares simulated method of moments with neural-network surrogate models using financial data from AT\&T and T-Mobile. This has trained me to formulate estimable mechanisms, select informative moments, assess uncertainty and sensitivity, and evaluate the predictive implications of simulation models. My principal expertise is complex-systems modelling rather than technology forecasting as a separate specialism. However, the combination of structural mechanisms, simulation, empirical calibration, and out-of-sample model assessment gives me a strong base from which to develop and test advanced forecasting methods within the programme.

\subsection*{Economics of innovation, technological change, and artificial intelligence}

My research examines technological investment and market evolution in ICT industries. The job-market paper studies strategic investment in service capacity and how it interacts with network effects, congestion, profitability, quality, and consumer adoption. My paper \textit{Heterogeneous Consumers and the Data Value Chain} analyses how data regulation changes firm incentives, market structure, prices, and consumer welfare across software markets. Together, these projects address core questions in the economics of technological change: how firms choose technological investments, how adoption depends on heterogeneous users and social interaction, and how institutions shape diffusion and market outcomes.

Artificial intelligence is a second active research stream. In the coauthored paper \textit{Agentic Workflows for Economic Research: Design and Implementation}, my collaborators and I develop a modular framework in which specialised LLM agents support tasks across the research lifecycle under human oversight. The implementation uses Microsoft's AutoGen and combines LLMs with tool use, multimodal inputs, workflow documentation, and reproducible research practices. This work gives me both conceptual knowledge of current AI capabilities and practical experience of designing research processes around them.

My previous work is not specialised in net-zero technologies. I would address this gap directly by combining systematic literature review with the programme's technology database, learning the relevant engineering and sectoral measures, and applying my existing expertise in investment, adoption, network effects, and diffusion to energy generation, transport, heat, industry, and carbon capture. These are precisely the settings in which feedback, heterogeneity, infrastructure constraints, and strategic investment make a complex-systems approach valuable.

\subsection*{Quantitative skills, programming, and sophisticated use of LLMs}

I bring a broad quantitative toolkit developed through economic theory, statistical analysis, simulation, calibration, and machine learning. I use Python for agent-based modelling, data processing, machine learning, and reproducible research workflows; Julia for computational economics and numerical methods; and R for statistics and econometrics. I also have experience with Mathematica, Stata, MATLAB, EViews, SQL, Linux, Git, simulated method of moments, MCMC, XGBoost, random forests, and neural networks.

My LLM work goes beyond conversational use. In the agentic-workflows project, I helped design a system of specialised agents that divides research tasks into explicit stages, connects models to tools and data, and retains human review, traceability, and version control. I have worked with AutoGen, LangChain, and the OpenAI API, and I use Claude Code in AI-assisted coding and research workflows. This combination would allow me to contribute to LLM-based data discovery, extraction, classification, validation, and documentation while treating provenance and quality control as core research requirements.

\subsection*{Contribution to peer-reviewed publications}

I contributed to \textit{Polytope Fraud Theory}, published in the \textit{International Review of Financial Analysis} in 2025. The study converts qualitative warning signs used by activist short sellers into ten measurable financial indicators and evaluates their ability to identify fraudulent firms using supervised learning, including XGBoost. The project required feature construction, analysis of heterogeneous financial information, comparison of model performance, and interpretation of statistical results in economic terms. I also have the coauthored LLM paper available as an arXiv preprint, two single-authored working papers, and the single-authored job-market paper from my dissertation. These projects demonstrate that I can develop a study from its conceptual and methodological foundations through analysis, writing, revision, and dissemination. I would bring that experience to the post's principal responsibility of leading and collaborating on peer-reviewed publications.

\subsection*{New research ideas and research-income generation}

Across my dissertation and collaborative work, I have developed research questions that connect economic mechanisms with new computational methods. Examples include estimating agent-based models with neural-network surrogates, studying how data restrictions alter software-market dynamics, and designing auditable LLM-agent workflows for economic research. For this programme, I would contribute ideas around three connected questions: whether technology-level characteristics improve forecasts beyond historical trend extrapolation; how uncertainty in cost forecasts propagates into adoption forecasts; and whether LLM-assisted extraction can expand technology datasets without weakening provenance or comparability.

My PhD was conducted in the EU-funded EPOC network, and my final stage was supported by a competitive Doctoral Completion Scholarship from the Bielefeld Young Researchers' Fund. This has given me experience of working within funded research programmes and connecting individual outputs to a wider collaborative agenda. I have not yet led an external grant application. I am ready to contribute literature synthesis, methodological design, work-package planning, preliminary analysis, and clear statements of expected contribution to proposals developed with senior researchers, while learning Oxford's research-development processes.

\subsection*{Communication, publication writing, and representation}

I communicate quantitative research through papers, conference presentations, workshops, and teaching materials. I have presented at the Workshop on Economics with Heterogeneous Interacting Agents, the World Congress of the International Microsimulation Association, the German Network for New Economic Dynamics, the Economic Policy in Complex Environments Conference, and other international meetings. These settings have required me to explain model assumptions, empirical strategies, and limitations to audiences from different disciplines.

I also co-organised and spoke at the public workshop ``AI Agents in Science'' at Wissenswerkstadt Bielefeld in 2025, presenting emerging research tools to a non-specialist audience in German. I served on the organising team for the International EPOC Doctoral Workshop at Ca' Foscari University of Venice and for the BiGSEM workshop series. Alongside my publication record, these experiences show that I can write for academic audiences, present results and proposals clearly, represent a research group professionally, and contribute to the wider intellectual life of the Smith School, INET, and the Oxford Martin School.

\section*{Desirable selection criteria}

\subsection*{Technology data: costs, diffusion, characteristics, and capabilities}

My existing empirical work offers relevant, although not identical, experience. For the parameter-estimation project, I organise firm-level financial data from AT\&T and T-Mobile to calibrate an ICT-industry model and compare estimation methods. My research on software markets represents adoption and information diffusion across heterogeneous consumers, while the job-market paper links firm investment and service capacity to quality, congestion, demand, and market outcomes. In \textit{Polytope Fraud Theory}, I helped convert unstructured qualitative allegations and public financial data into a manually labelled dataset and economically interpretable features.

I have not yet assembled a cross-technology dataset of historical costs and capabilities. I would approach that responsibility by first establishing a transparent data model, common definitions, units, sources, and versioning rules; then building reproducible Python workflows for collection, validation, harmonisation, and documentation; and finally testing sensitivity to missing observations, revisions, and measurement choices. My experience moving between economic concepts, heterogeneous source material, structured features, and computational analysis is well suited to this work.

\subsection*{Theory and empirics of technological change and innovation}

My strongest substantive preparation lies in the economics of ICT investment, adoption, network effects, data governance, and AI. This is complemented by empirical calibration and machine-learning research. I can therefore contribute immediately to the programme's work on artificial intelligence and to general questions about technological characteristics, firm incentives, diffusion, and feedback. I would extend this foundation to net-zero technologies through close collaboration with domain specialists and systematic engagement with evidence on learning curves, deployment, performance, and sector-specific constraints. The opportunity to make comparisons across AI and clean technologies is a central reason I want to join this programme.

\section*{Research management, collaboration, and motivation}

I have managed several single-authored projects from question formulation through modelling, data work, analysis, and writing, while also contributing to multi-institutional projects with coauthors from economics, finance, and artificial intelligence. Working across Bielefeld University, Universit\`a Cattolica del Sacro Cuore, and the EPOC network has required structured planning, transparent documentation, constructive feedback, and reliable delivery across institutions. Workshop organisation has added practical experience of coordinating contributors and deadlines.

I have not yet formally supervised a team of research assistants. I would bring to that responsibility a concrete management approach: divide the database into well-defined work packages, provide written protocols and examples, use version-controlled task and data documentation, review samples early, and hold short regular meetings to resolve ambiguities before they propagate. I would combine clear quality standards with space for research assistants to raise problems and suggest improvements. This approach follows directly from my experience building reproducible research and agentic workflows in which responsibilities, inputs, outputs, and checks must be explicit.

I want to undertake this role because it joins three questions that define my research trajectory: how technologies improve, how firm and user decisions shape their diffusion, and how quantitative methods can produce credible forecasts in complex systems. The combination of the Complexity Economics research group, the Oxford Martin Programme, and collaboration across disciplines and partner institutions would allow me to contribute my established strengths in computational economics and AI while developing rigorous comparative research on technological progress. I would aim to deliver well-documented data, transparent and testable methods, strong peer-reviewed publications, and constructive support for the wider research group.

\end{document}
